\subsubsection{Mixed surface boundaries with finite exceptional wild image}

\begin{theorem}[Proved under stated hypotheses: mixed finite boundary with finite cone covers]
\label{D:finite-cone-mixed-boundary}
Let $k$ be algebraically closed of characteristic $p>0$, let
$S$ be a normal connected proper surface over $k$, and let
$Y\subset S$ be finite.  Suppose every singular point of $Y$
satisfies at least one of the following conditions:
\begin{enumerate}
\item $S$ is $\mathbb Q$-factorial at that point;
\item $S$ is log canonical there;
\item its completed local ring is a completed positively graded
      normal surface cone;
\item a projective resolution of its germ has sole reduced
      exceptional fibre a smooth proper curve $C$, and a finite
      map $h:C'\to C$ satisfies the three conditions
      \eqref{D:cone-first}--\eqref{D:cone-normal}, or some
      $C_r=(C,\mathcal O_{\widetilde S}/I_C^r)$ retracts onto
      $C$ with $r\geq2$ and $r(-C^2)>2g(C)-2$;
\item its completed local ring is isomorphic to one of the
      explicit genus-two completed germs $S$ or $T$ of
      Example~\ref{D:wild-genus2-single-curve};
\item a projective resolution has sole reduced exceptional
      fibre a smooth curve of negative self-intersection, and
      its entire formal neighborhood admits a retraction to
      that curve.
\end{enumerate}
Then $\pi_1(S\setminus Y)$ is topologically finitely generated.
The singularities retained in the open are unrestricted.

In (4), if $C^2=-d<0$ and $g=g(C)$, either of the following
conditions suffices:
\[
 d>\max\{0,2g-2\};\qquad
 g\geq2,\quad d=2g-2,\quad
 \omega_C\otimes N_{C/\widetilde S}\text{ is torsion}.
\]
In the equality case, the torsion condition holds whenever
the germ is $\mathbb Q$-Gorenstein, with arbitrary Cartier index.
If the first neighborhood $C_2$ retracts onto $C$, the weaker
numerical condition $d>\max\{0,g-1\}$ suffices.

The same conclusion holds for a possibly nonnormal proper
scheme of dimension at most two when $Y$ has codimension at
least two in every irreducible component and its inverse image
on every normal surface component satisfies these alternatives.
\end{theorem}
\begin{proof}
Corollaries~\ref{D:negative-curve-finite-wild}
and~\ref{D:negative-curve-equality-finite-wild} prove the stated
sufficient conditions for (4); Proposition~\ref{D:finite-jet-cone-cover}
and Corollary~\ref{D:split-first-neighborhood-finite-wild} supply
its finite-order retraction alternatives.  We prove the general assertion
with alternatives (1)--(6).

We construct a normal proper $\mathbb Q$-factorial modification
$q:X\to S$.  Leave the omitted $\mathbb Q$-factorial singular
points unchanged.  At every other omitted point of type (2),
use its minimal resolution and, for a nodal cyclic core, the
necessary node subdivisions.  A rational-tree exceptional
configuration is already $\mathbb Q$-factorial, by
Proposition~\ref{D:rational-tree-Qfactorial}, so that case was
retained.  At points of type (3), use the partial extraction
of Proposition~\ref{D:weighted-graded-extraction}; it is
$\mathbb Q$-factorial and has just its radial exceptional curve.
At points of type (4), use the specified resolution.  At points of type (6), use the resolution carrying the formal
retraction.  At points
of type (5), the explicit one-curve resolution descends from
completion.  Indeed, use the image weighted Rees algebra
$\bigoplus_m (I_m\mathcal O_S)s^m\subset\mathcal O_S[s]$
of the explicit strict transform, not the whole ambient
base-change blowup.  A sufficiently divisible Veronese is
generated in degree one and equals the Rees algebra of
$J=I_N\mathcal O_S$.  Its Proj is therefore the blowup of
that ideal, which is primary to the closed point; see \cite[Lemma 10.56.2, tag 0EGH]{C_StacksWeightedCone}.
Such an ideal descends from the completion since it contains
a power of the maximal ideal.  Its blowup therefore descends
with the same completed model and the same exceptional curve,
by \cite[Lemmas 54.11.1--3, section 0BG4]{StacksBFiniteCone}.
Regularity descends faithfully flatly.  In particular we do
not replace this one-curve model by a dominating blowup with
additional, uncontrolled exceptional divisors.
Finally resolve all singular points retained in $S\setminus Y$.

There are only finitely many singular points.  The local
modifications are isomorphisms on punctured neighborhoods;
spreading and gluing with the identity therefore give the
claimed proper birational map.  Its source is regular except
at the retained or extracted $\mathbb Q$-factorial points.
It is thus $\mathbb Q$-factorial everywhere and projective by
Lemma~\ref{D:Qfactorial-partial-resolution}.

Put $U=S\setminus Y$ and $U'=q^{-1}(U)$.  Connected finite
\'etale covers of the normal integral $U$ stay connected on
$U'$, since their generic algebra is unchanged.  Consequently
$\pi_1(U')\to\pi_1(U)$ is surjective.  Every exceptional
divisor of $X\setminus U'$ lies over an omitted point of one
of the types just resolved or partially extracted.

For type (2), its wild inertia image in $\pi_1(U)$ is zero by
Corollary~\ref{D:lc-core-valuative-tame}.  For type (3), it is
zero by Proposition~\ref{D:graded-cone-tame} and the weighted
extraction comparison.  For type (4), it is finite by
Proposition~\ref{D:finite-cone-wild-bound} or
Proposition~\ref{D:finite-jet-cone-cover}.  The same is true
for type (5), with bounds four or eight, by
Corollary~\ref{D:wild-genus2-finite-inertia}.
For type (6), the wild image is zero by
Proposition~\ref{D:formal-retraction-tame}.  These are bounds
in the full punctured local group, so they bound the images
in the global group as well.  The completion comparison uses
the same completed divisorial DVR for every original cover.
Retained omitted points have codimension two on $X$ and create
no divisorial boundary.  Exceptional curves above singularities
retained in $U$ lie in $U'$.

We may therefore apply Proposition~\ref{D:finite-exceptional-wild}
to this projective modification.  It chooses one open normal
subgroup avoiding the finitely many finite wild images; after
normalization in its finite cover, the resulting open subgroup
is a quotient of a projective fixed-boundary tame group.
That subgroup, and hence $\pi_1(U)$, is finitely generated.
In particular we have used finiteness of these wild images,
without asserting that the original covers are tame at them.

For a nonnormal proper scheme, pass to the disjoint finite
normalizations of its reduced irreducible components.  The
normal surface cases have just been proved.  Components of
dimension at most one have no omitted points under the
codimension hypothesis, and their proper fundamental groups
are finitely generated.  Finite proper descent, in the form
used in Proposition~\ref{D:fixed-boundary-positive} and
\cite[IX, Corollary 5.2]{SGA1}, proves the assertion downstairs.
\end{proof}

\begin{corollary}[Proved under stated hypotheses: isolated fixed boundary with finite cone covers]
\label{D:finite-cone-fixed-boundary}
Let $S$ be a normal connected proper surface over an
algebraically closed field of positive characteristic, and
let $Y\subsetneq S$ be closed.  Suppose every singular point
of $Y$ which is not $\mathbb Q$-factorial is isolated in $Y$
and satisfies one of alternatives (2)--(6) of
Theorem~\ref{D:finite-cone-mixed-boundary}.  Then $Q(S,Y)$ is
topologically finitely generated.  The numerical and torsion
sufficient conditions stated there apply, and retained
singularities are unrestricted.
\end{corollary}
\begin{proof}
Make the same partial projective modification as in the theorem,
retaining all $\mathbb Q$-factorial boundary singularities.
The strict transforms of original divisorial boundary points
have the identical DVRs, so their wild inertia already has
zero image in $Q(S,Y)$.  Every additional boundary divisor
lies over one of the specified isolated points.  Isolation
ensures that every cover counted by $Q(S,Y)$ restricts to a
finite \'etale cover of that punctured germ.  The preceding
local results consequently give a finite wild image at each
such divisor.  Curves introduced over points outside $Y$
remain in the open.  Proposition~\ref{D:finite-exceptional-wild}
again applies and proves finite generation.  This argument
retains the isolation hypothesis at the newly resolved
boundary points.
\end{proof}
